\subsection{Lazy regional selection from training inputs}\label{sec:lazy_regional}

The regional bound can drive computation before a full continuation partition exists. The following implementation takes the training responses, fixed parameters, first-action menu and eligible design pools. The computation discovers its action domains and constructs its conditional posteriors from these inputs.

After constructing the raw-mean coefficient bands and enclosing the first-response marginals, initialize a response-cell cover on $[-8,8]$, split at the enclosed incumbent crossing. Each unresolved cell uses $[B_C^-,m_0 P_C^+]$, where $B_C$ is its first-evaluation improvement and $P_C$ its probability. A uniform-reward control uses $[0,m_0P_C^+]$ instead. The explicit outside-window tail is retained.

Choose the largest-width cell among the current best-lower-bound branch and its highest-upper-bound competitor. If the cell has no certified action, evaluate the nominal affine selector at its center and perform one linear dominance scan for that reference. Intersect the resulting stability domain with the response cell using outward parameter transformations. Assign the certified intersection and retain unresolved remainder cells. Only an encountered action requires a bivariate conditional posterior. On a certified cell, the fixed-action conditional integral and its curvature remainder are used; subdividing that cell reduces integration error. The process ends with the selection inequality or an unresolved outcome at its resource limit. No global affine-envelope construction is used.

\paragraph{Guarantee with uncomputed conditional laws.}
For each first action, let $\mathcal P_a$ contain the positive conditional response laws that agree with every queried conditional marginal on its certified region. An unqueried response is restricted only by positivity. Each unresolved cell then has a valid reward range, and every resolved cell has its checked conditional integral. Thus, if $[\underline V_a,\overline V_a]$ is the assembled interval,
\begin{equation}
\sup_{P_a\in\mathcal P_a}V_a(P_a)\le\overline V_a,
\qquad
\inf_{P_a\in\mathcal P_a}V_a(P_a)\ge\underline V_a.
\label{eq:lazy_completion}
\end{equation}
The selection rule therefore bounds first-action model regret for every admissible positive completion of the uncomputed laws. The uncertainty class represents the uncomputed part of the planning expectation.

Each queried bivariate posterior is constructed from the same fixed binary training matrix and normalization as the earlier GP specification, with residual and positive-variance checks. The lazy implementation does not assemble or certify the full unqueried joint covariance. If a joint GP with these conditional marginals exists, its rollout value is included. Positive conditional laws also define the branch expectations without requiring correlations between unrealized counterfactual second actions. 

\paragraph{Local error and computational scope.}
With exact affine coefficients and exact model parameters, let $J_a$ be the number of positive-length policy switches within $[-T,T]$, and let $C_a$ uniformly bound the smooth conditional-reward curvature. Consider a local cell cover of maximum width $h$ after one center probe per cell. Cells containing no switch have certified identities; the remaining cells have total length at most $J_a h$. The resulting value-interval width is at most
\begin{equation}
\frac{m_0 J_a h}{\sqrt{2\pi}}+
\frac{C_a h^2}{4}+2m_0\Phi(-T).
\label{eq:lazy_local_error}
\end{equation}
The first term covers unresolved switch cells, the second the two-sided interpolation remainder, and the third the tail. An affine lower envelope has at most one positive-length region per distinct line, so $J_a$ is finite. This idealized bound explains how local refinement handles discontinuities without a global policy-value Lipschitz constant. Fixed parameter bands and finite precision can retain an uncertainty floor, and the finite-budget schedule can return unresolved.

With $S$ local discovery scans over $N$ eligible continuations, dominance work is $O(SN)$ before exact-arithmetic bit costs; posterior and integration costs are additional. The decision inequality determines whether limited regions and posterior requests suffice. The minimal control changes only the unresolved reward range, so its comparison isolates the effect of the regional bound within the same lazy discovery algorithm.
